\documentclass[11pt,a4paper]{article}

% -----------------------------------------------------------------------------
% Geometry and base packages
% -----------------------------------------------------------------------------
\usepackage[margin=1in]{geometry}
\usepackage[T1]{fontenc}
\usepackage[utf8]{inputenc}

% Prefer scalable fonts when available (microtype expansion); fall back silently.
\IfFileExists{lmodern.sty}{\usepackage{lmodern}}{}

% Shared macros for PYNE arXiv paper series
% Include with: % Shared macros for PYNE arXiv paper series
% Include with: % Shared macros for PYNE arXiv paper series
% Include with: \input{../common/macros.tex}

\usepackage{amsmath,amssymb,amsfonts}
\usepackage{amsthm}
\usepackage{booktabs}
\usepackage{graphicx}
\usepackage{xcolor}
\usepackage{hyperref}
\usepackage{url}
\usepackage{algorithm}
\usepackage{algpseudocode}
\usepackage{listings}
\usepackage{tikz}
\usepackage{pgfplots}
\pgfplotsset{compat=1.17}
\usetikzlibrary{arrows.meta,positioning,fit,calc,shapes.geometric,shapes.multipart,decorations.pathreplacing,backgrounds,chains,matrix,patterns}
\usepackage{subcaption}
\usepackage{multirow}
\usepackage{tabularx}
\usepackage{enumitem}
\usepackage{microtype}
\usepackage{balance}
\usepackage{csquotes}
\usepackage{natbib}

% Colors
\definecolor{pyneorange}{HTML}{F97316}
\definecolor{pyneblue}{HTML}{2563EB}
\definecolor{pynegray}{HTML}{6B7280}
\definecolor{pynegreen}{HTML}{16A34A}
\definecolor{codebg}{HTML}{F8FAFC}
\definecolor{codekw}{HTML}{7C3AED}
\definecolor{codestr}{HTML}{B45309}
\definecolor{codecomment}{HTML}{64748B}

\hypersetup{
  colorlinks=true,
  linkcolor=pyneblue,
  citecolor=pynegreen,
  urlcolor=pyneorange,
  pdftitle={},
  pdfauthor={HOOX / PYNE Project}
}

% Code listings
\lstdefinestyle{pine}{
  basicstyle=\ttfamily\footnotesize,
  backgroundcolor=\color{codebg},
  keywordstyle=\color{codekw}\bfseries,
  stringstyle=\color{codestr},
  commentstyle=\color{codecomment}\itshape,
  breaklines=true,
  frame=single,
  rulecolor=\color{black!20},
  numbers=left,
  numberstyle=\tiny\color{pynegray},
  numbersep=6pt,
  xleftmargin=1.5em,
  framexleftmargin=1.2em,
  showstringspaces=false,
  tabsize=2,
  morekeywords={indicator,strategy,library,plot,plotshape,fill,hline,var,varip,if,else,for,while,switch,true,false,na,and,or,not,import,export,type,method,enum},
  morecomment=[l]{//},
  morestring=[b]",
}
\lstdefinestyle{python}{
  language=Python,
  basicstyle=\ttfamily\footnotesize,
  backgroundcolor=\color{codebg},
  keywordstyle=\color{codekw}\bfseries,
  stringstyle=\color{codestr},
  commentstyle=\color{codecomment}\itshape,
  breaklines=true,
  frame=single,
  rulecolor=\color{black!20},
  numbers=left,
  numberstyle=\tiny\color{pynegray},
  numbersep=6pt,
  xleftmargin=1.5em,
  framexleftmargin=1.2em,
  showstringspaces=false,
  tabsize=4,
}
\lstset{style=python}

% Theorem environments
\theoremstyle{plain}
\newtheorem{theorem}{Theorem}[section]
\newtheorem{lemma}[theorem]{Lemma}
\newtheorem{proposition}[theorem]{Proposition}
\newtheorem{corollary}[theorem]{Corollary}
\theoremstyle{definition}
\newtheorem{definition}[theorem]{Definition}
\newtheorem{example}[theorem]{Example}
\newtheorem{remark}[theorem]{Remark}

% Shortcuts
\newcommand{\pyne}{\textsc{Pyne}}
\newcommand{\pine}{Pine Script\texttrademark}
\newcommand{\tv}{TradingView\textregistered}
\newcommand{\asdl}{\textsc{Asdl}}
\newcommand{\antlr}{\textsc{Antlr}4}
\newcommand{\numba}{\textsc{Numba}}
\newcommand{\lsp}{\textsc{Lsp}}
\newcommand{\ohlcv}{\textsc{Ohlcv}}
\newcommand{\udt}{\textsc{Udt}}
\newcommand{\jit}{\textsc{Jit}}
\newcommand{\api}{\textsc{Api}}
\newcommand{\cli}{\textsc{Cli}}
% AST acronym (text mode). Keep amsmath's math asterisk \ast intact.
\DeclareRobustCommand{\Ast}{\textsc{Ast}}
% Papers in this series write \ast{} for the AST acronym in prose.
% Redefine robustly; use \mathast for a binary asterisk in formulas.
\let\mathast\ast
\DeclareRobustCommand{\ast}{\textsc{Ast}}
\newcommand{\ir}{\textsc{Ir}}
\newcommand{\ta}{\texttt{ta}}
\newcommand{\code}[1]{\texttt{#1}}
\newcommand{\file}[1]{\texttt{#1}}
\newcommand{\eg}{e.g.}
\newcommand{\ie}{i.e.}
\newcommand{\etal}{\textit{et~al.}}
\newcommand{\resp}{respectively}

% Math helpers
\newcommand{\R}{\mathbb{R}}
\newcommand{\N}{\mathbb{N}}
\newcommand{\Z}{\mathbb{Z}}
% Text-mode NA token for prose; use $\namath$ inside formulas if needed.
\DeclareRobustCommand{\na}{\textsf{na}}
\newcommand{\namath}{\mathsf{na}}
\newcommand{\series}[1]{\mathbf{#1}}
\newcommand{\baridx}{i}
\newcommand{\bars}{n}
\newcommand{\period}{p}
\newcommand{\abserr}{\varepsilon_{\mathrm{abs}}}
\newcommand{\relerr}{\varepsilon_{\mathrm{rel}}}
\newcommand{\maxerr}{\varepsilon_{\mathrm{max}}}
\newcommand{\meanerr}{\varepsilon_{\mathrm{mean}}}

% Pipeline notation — use inside math: $\pipeline$
\newcommand{\pipeline}{\mathrm{src}\!\to\!\mathrm{lex}\!\to\!\mathrm{parse}\!\to\!\mathrm{AST}\!\to\!\mathrm{eval}}

% Disclaimer footnote helper
\newcommand{\trademarknote}{%
  \pine{} and \tv{} are trademarks of TradingView, Inc.
  \pyne{} is an independent, unofficial implementation and is not affiliated with,
  authorized by, sponsored by, or endorsed by TradingView, Inc.
}

% Figure path helper (papers live one level under .papers/)
\graphicspath{{figures/}{../common/figures/}}


\usepackage{amsmath,amssymb,amsfonts}
\usepackage{amsthm}
\usepackage{booktabs}
\usepackage{graphicx}
\usepackage{xcolor}
\usepackage{hyperref}
\usepackage{url}
\usepackage{algorithm}
\usepackage{algpseudocode}
\usepackage{listings}
\usepackage{tikz}
\usepackage{pgfplots}
\pgfplotsset{compat=1.17}
\usetikzlibrary{arrows.meta,positioning,fit,calc,shapes.geometric,shapes.multipart,decorations.pathreplacing,backgrounds,chains,matrix,patterns}
\usepackage{subcaption}
\usepackage{multirow}
\usepackage{tabularx}
\usepackage{enumitem}
\usepackage{microtype}
\usepackage{balance}
\usepackage{csquotes}
\usepackage{natbib}

% Colors
\definecolor{pyneorange}{HTML}{F97316}
\definecolor{pyneblue}{HTML}{2563EB}
\definecolor{pynegray}{HTML}{6B7280}
\definecolor{pynegreen}{HTML}{16A34A}
\definecolor{codebg}{HTML}{F8FAFC}
\definecolor{codekw}{HTML}{7C3AED}
\definecolor{codestr}{HTML}{B45309}
\definecolor{codecomment}{HTML}{64748B}

\hypersetup{
  colorlinks=true,
  linkcolor=pyneblue,
  citecolor=pynegreen,
  urlcolor=pyneorange,
  pdftitle={},
  pdfauthor={HOOX / PYNE Project}
}

% Code listings
\lstdefinestyle{pine}{
  basicstyle=\ttfamily\footnotesize,
  backgroundcolor=\color{codebg},
  keywordstyle=\color{codekw}\bfseries,
  stringstyle=\color{codestr},
  commentstyle=\color{codecomment}\itshape,
  breaklines=true,
  frame=single,
  rulecolor=\color{black!20},
  numbers=left,
  numberstyle=\tiny\color{pynegray},
  numbersep=6pt,
  xleftmargin=1.5em,
  framexleftmargin=1.2em,
  showstringspaces=false,
  tabsize=2,
  morekeywords={indicator,strategy,library,plot,plotshape,fill,hline,var,varip,if,else,for,while,switch,true,false,na,and,or,not,import,export,type,method,enum},
  morecomment=[l]{//},
  morestring=[b]",
}
\lstdefinestyle{python}{
  language=Python,
  basicstyle=\ttfamily\footnotesize,
  backgroundcolor=\color{codebg},
  keywordstyle=\color{codekw}\bfseries,
  stringstyle=\color{codestr},
  commentstyle=\color{codecomment}\itshape,
  breaklines=true,
  frame=single,
  rulecolor=\color{black!20},
  numbers=left,
  numberstyle=\tiny\color{pynegray},
  numbersep=6pt,
  xleftmargin=1.5em,
  framexleftmargin=1.2em,
  showstringspaces=false,
  tabsize=4,
}
\lstset{style=python}

% Theorem environments
\theoremstyle{plain}
\newtheorem{theorem}{Theorem}[section]
\newtheorem{lemma}[theorem]{Lemma}
\newtheorem{proposition}[theorem]{Proposition}
\newtheorem{corollary}[theorem]{Corollary}
\theoremstyle{definition}
\newtheorem{definition}[theorem]{Definition}
\newtheorem{example}[theorem]{Example}
\newtheorem{remark}[theorem]{Remark}

% Shortcuts
\newcommand{\pyne}{\textsc{Pyne}}
\newcommand{\pine}{Pine Script\texttrademark}
\newcommand{\tv}{TradingView\textregistered}
\newcommand{\asdl}{\textsc{Asdl}}
\newcommand{\antlr}{\textsc{Antlr}4}
\newcommand{\numba}{\textsc{Numba}}
\newcommand{\lsp}{\textsc{Lsp}}
\newcommand{\ohlcv}{\textsc{Ohlcv}}
\newcommand{\udt}{\textsc{Udt}}
\newcommand{\jit}{\textsc{Jit}}
\newcommand{\api}{\textsc{Api}}
\newcommand{\cli}{\textsc{Cli}}
% AST acronym (text mode). Keep amsmath's math asterisk \ast intact.
\DeclareRobustCommand{\Ast}{\textsc{Ast}}
% Papers in this series write \ast{} for the AST acronym in prose.
% Redefine robustly; use \mathast for a binary asterisk in formulas.
\let\mathast\ast
\DeclareRobustCommand{\ast}{\textsc{Ast}}
\newcommand{\ir}{\textsc{Ir}}
\newcommand{\ta}{\texttt{ta}}
\newcommand{\code}[1]{\texttt{#1}}
\newcommand{\file}[1]{\texttt{#1}}
\newcommand{\eg}{e.g.}
\newcommand{\ie}{i.e.}
\newcommand{\etal}{\textit{et~al.}}
\newcommand{\resp}{respectively}

% Math helpers
\newcommand{\R}{\mathbb{R}}
\newcommand{\N}{\mathbb{N}}
\newcommand{\Z}{\mathbb{Z}}
% Text-mode NA token for prose; use $\namath$ inside formulas if needed.
\DeclareRobustCommand{\na}{\textsf{na}}
\newcommand{\namath}{\mathsf{na}}
\newcommand{\series}[1]{\mathbf{#1}}
\newcommand{\baridx}{i}
\newcommand{\bars}{n}
\newcommand{\period}{p}
\newcommand{\abserr}{\varepsilon_{\mathrm{abs}}}
\newcommand{\relerr}{\varepsilon_{\mathrm{rel}}}
\newcommand{\maxerr}{\varepsilon_{\mathrm{max}}}
\newcommand{\meanerr}{\varepsilon_{\mathrm{mean}}}

% Pipeline notation — use inside math: $\pipeline$
\newcommand{\pipeline}{\mathrm{src}\!\to\!\mathrm{lex}\!\to\!\mathrm{parse}\!\to\!\mathrm{AST}\!\to\!\mathrm{eval}}

% Disclaimer footnote helper
\newcommand{\trademarknote}{%
  \pine{} and \tv{} are trademarks of TradingView, Inc.
  \pyne{} is an independent, unofficial implementation and is not affiliated with,
  authorized by, sponsored by, or endorsed by TradingView, Inc.
}

% Figure path helper (papers live one level under .papers/)
\graphicspath{{figures/}{../common/figures/}}


\usepackage{amsmath,amssymb,amsfonts}
\usepackage{amsthm}
\usepackage{booktabs}
\usepackage{graphicx}
\usepackage{xcolor}
\usepackage{hyperref}
\usepackage{url}
\usepackage{algorithm}
\usepackage{algpseudocode}
\usepackage{listings}
\usepackage{tikz}
\usepackage{pgfplots}
\pgfplotsset{compat=1.17}
\usetikzlibrary{arrows.meta,positioning,fit,calc,shapes.geometric,shapes.multipart,decorations.pathreplacing,backgrounds,chains,matrix,patterns}
\usepackage{subcaption}
\usepackage{multirow}
\usepackage{tabularx}
\usepackage{enumitem}
\usepackage{microtype}
\usepackage{balance}
\usepackage{csquotes}
\usepackage{natbib}

% Colors
\definecolor{pyneorange}{HTML}{F97316}
\definecolor{pyneblue}{HTML}{2563EB}
\definecolor{pynegray}{HTML}{6B7280}
\definecolor{pynegreen}{HTML}{16A34A}
\definecolor{codebg}{HTML}{F8FAFC}
\definecolor{codekw}{HTML}{7C3AED}
\definecolor{codestr}{HTML}{B45309}
\definecolor{codecomment}{HTML}{64748B}

\hypersetup{
  colorlinks=true,
  linkcolor=pyneblue,
  citecolor=pynegreen,
  urlcolor=pyneorange,
  pdftitle={},
  pdfauthor={HOOX / PYNE Project}
}

% Code listings
\lstdefinestyle{pine}{
  basicstyle=\ttfamily\footnotesize,
  backgroundcolor=\color{codebg},
  keywordstyle=\color{codekw}\bfseries,
  stringstyle=\color{codestr},
  commentstyle=\color{codecomment}\itshape,
  breaklines=true,
  frame=single,
  rulecolor=\color{black!20},
  numbers=left,
  numberstyle=\tiny\color{pynegray},
  numbersep=6pt,
  xleftmargin=1.5em,
  framexleftmargin=1.2em,
  showstringspaces=false,
  tabsize=2,
  morekeywords={indicator,strategy,library,plot,plotshape,fill,hline,var,varip,if,else,for,while,switch,true,false,na,and,or,not,import,export,type,method,enum},
  morecomment=[l]{//},
  morestring=[b]",
}
\lstdefinestyle{python}{
  language=Python,
  basicstyle=\ttfamily\footnotesize,
  backgroundcolor=\color{codebg},
  keywordstyle=\color{codekw}\bfseries,
  stringstyle=\color{codestr},
  commentstyle=\color{codecomment}\itshape,
  breaklines=true,
  frame=single,
  rulecolor=\color{black!20},
  numbers=left,
  numberstyle=\tiny\color{pynegray},
  numbersep=6pt,
  xleftmargin=1.5em,
  framexleftmargin=1.2em,
  showstringspaces=false,
  tabsize=4,
}
\lstset{style=python}

% Theorem environments
\theoremstyle{plain}
\newtheorem{theorem}{Theorem}[section]
\newtheorem{lemma}[theorem]{Lemma}
\newtheorem{proposition}[theorem]{Proposition}
\newtheorem{corollary}[theorem]{Corollary}
\theoremstyle{definition}
\newtheorem{definition}[theorem]{Definition}
\newtheorem{example}[theorem]{Example}
\newtheorem{remark}[theorem]{Remark}

% Shortcuts
\newcommand{\pyne}{\textsc{Pyne}}
\newcommand{\pine}{Pine Script\texttrademark}
\newcommand{\tv}{TradingView\textregistered}
\newcommand{\asdl}{\textsc{Asdl}}
\newcommand{\antlr}{\textsc{Antlr}4}
\newcommand{\numba}{\textsc{Numba}}
\newcommand{\lsp}{\textsc{Lsp}}
\newcommand{\ohlcv}{\textsc{Ohlcv}}
\newcommand{\udt}{\textsc{Udt}}
\newcommand{\jit}{\textsc{Jit}}
\newcommand{\api}{\textsc{Api}}
\newcommand{\cli}{\textsc{Cli}}
% AST acronym (text mode). Keep amsmath's math asterisk \ast intact.
\DeclareRobustCommand{\Ast}{\textsc{Ast}}
% Papers in this series write \ast{} for the AST acronym in prose.
% Redefine robustly; use \mathast for a binary asterisk in formulas.
\let\mathast\ast
\DeclareRobustCommand{\ast}{\textsc{Ast}}
\newcommand{\ir}{\textsc{Ir}}
\newcommand{\ta}{\texttt{ta}}
\newcommand{\code}[1]{\texttt{#1}}
\newcommand{\file}[1]{\texttt{#1}}
\newcommand{\eg}{e.g.}
\newcommand{\ie}{i.e.}
\newcommand{\etal}{\textit{et~al.}}
\newcommand{\resp}{respectively}

% Math helpers
\newcommand{\R}{\mathbb{R}}
\newcommand{\N}{\mathbb{N}}
\newcommand{\Z}{\mathbb{Z}}
% Text-mode NA token for prose; use $\namath$ inside formulas if needed.
\DeclareRobustCommand{\na}{\textsf{na}}
\newcommand{\namath}{\mathsf{na}}
\newcommand{\series}[1]{\mathbf{#1}}
\newcommand{\baridx}{i}
\newcommand{\bars}{n}
\newcommand{\period}{p}
\newcommand{\abserr}{\varepsilon_{\mathrm{abs}}}
\newcommand{\relerr}{\varepsilon_{\mathrm{rel}}}
\newcommand{\maxerr}{\varepsilon_{\mathrm{max}}}
\newcommand{\meanerr}{\varepsilon_{\mathrm{mean}}}

% Pipeline notation — use inside math: $\pipeline$
\newcommand{\pipeline}{\mathrm{src}\!\to\!\mathrm{lex}\!\to\!\mathrm{parse}\!\to\!\mathrm{AST}\!\to\!\mathrm{eval}}

% Disclaimer footnote helper
\newcommand{\trademarknote}{%
  \pine{} and \tv{} are trademarks of TradingView, Inc.
  \pyne{} is an independent, unofficial implementation and is not affiliated with,
  authorized by, sponsored by, or endorsed by TradingView, Inc.
}

% Figure path helper (papers live one level under .papers/)
\graphicspath{{figures/}{../common/figures/}}


% If only bitmap CM fonts are installed, disable microtype expansion to avoid
% "auto expansion is only possible with scalable fonts" fatal errors.
\makeatletter
\@ifpackageloaded{microtype}{%
  \microtypesetup{expansion=false}%
}{}
\makeatother

\hypersetup{
  pdftitle={Language Tooling for Financial DSLs: An LSP Architecture for Pine Script with Multi-Editor Delivery},
  pdfauthor={HOOX Research / PYNE Contributors},
  pdfkeywords={Language Server Protocol, developer tools, code completion, diagnostics, domain-specific languages, VS Code extensions, Pine Script}
}

\title{%
  Language Tooling for Financial DSLs:\\
  An \lsp{} Architecture for \pine{} with Multi-Editor Delivery}
\author{%
  HOOX Research / \pyne{} Contributors\\
  \texttt{jango\_blockchained}\\
  Independent Open Source Research\\
  \texttt{contact@hoox.sh}}
\date{}

\begin{document}
\maketitle

\begin{abstract}
Financial domain-specific languages (\textsc{Dsl}s) used for technical analysis and
algorithmic trading are often edited exclusively inside proprietary, browser-hosted
environments. That coupling denies practitioners the offline diagnostics, completion,
navigation, and multi-editor workflows that general-purpose language tooling takes for
granted. This paper presents the design and implementation of \texttt{pynescript-lsp},
the Language Server Protocol (\lsp{}) implementation for \pyne{}---an independent open
toolchain for \pine{}~\citep{pynescript2026,tradingview-pine}. The server is built on
\texttt{pygls}~\citep{murphy2020pygls}, shares a single \ast{} pipeline with the project's
parser, linter, unparser, and evaluator, and can be shipped either as a Python module
(\texttt{pip install hoox-pyne[lsp]}) or as a Nuitka onefile binary~\citep{nuitka}.

We describe a feature surface covering push and pull diagnostics, metadata-driven
completion over hundreds of builtins (\(\sim\)800+ catalog entries), hover documentation,
definition/references, document and workspace symbols, semantic tokens, inlay type
hints, and unparse-based formatting. A generation--encryption--integrity pipeline
protects the builtin documentation corpus in compiled distributions without coupling
editor features to network services. Thin clients for VS~Code, Neovim, Zed, and Emacs,
together with a complementary \cli{}, demonstrate multi-surface delivery. We report
implementation lessons for \lsp{} servers over financial \textsc{Dsl}s, discuss
boundaries with evaluation and chart/backtest \api{} surfaces, and situate the work in
the broader language-tooling ecosystem~\citep{lsp2023,voutilainen2016lsp,fowler2010dsl}.
\end{abstract}

\vspace{0.5em}
\noindent\textbf{Keywords:}
Language Server Protocol, developer tools, code completion, diagnostics,
domain-specific languages, VS~Code extensions, Pine Script

\vspace{0.75em}
\noindent\footnotesize\textit{\trademarknote}
\normalsize

\tableofcontents
\newpage

% =============================================================================
\section{Introduction}
\label{sec:intro}
% =============================================================================

Domain-specific languages for quantitative trading occupy an awkward middle ground
between spreadsheets and full programming environments. \pine{}, the scripting language
used by millions of traders on the \tv{} platform, exemplifies this tension: it is
expressive enough to encode multi-timeframe strategies, user-defined types (\udt{}s),
and rich drawing objects, yet its primary authoring surface remains a browser-only
editor tightly coupled to a hosted runtime~\citep{tradingview-pine,fowler2010dsl}.
Practitioners who prefer local Git workflows, keyboard-driven editors, continuous
integration, or offline research therefore lack the language services that have become
standard elsewhere: incremental diagnostics, ranked completion, hover documentation,
go-to-definition, and deterministic formatting~\citep{lsp2023,voutilainen2016lsp}.

The Language Server Protocol (\lsp{})~\citep{lsp2023} was introduced precisely to
decouple language intelligence from editor implementations. A single server process
exposes a JSON-RPC surface; any compliant client---VS~Code, Neovim, Zed, Emacs, Helix,
and others---can consume it. For general-purpose languages the ecosystem is mature
(TypeScript, Rust Analyzer, Pyright). For financial \textsc{Dsl}s the picture is
sparse: tooling is either proprietary and online, or limited to syntax highlighting
without semantic analysis.

\pyne{} (\texttt{hoox-pyne} on PyPI; import package \texttt{pynescript}) is an
independent, unofficial open toolchain for \pine{}~\citep{pynescript2026,hoox2026}.
It provides a lossless ANTLR4-based parser~\citep{parr2013antlr}, an ASDL-typed
\ast{}~\citep{wang1997asdl}, a static linter, an unparser, an interpreter/compiler
runtime, and---the subject of this paper---an \lsp{} server
(\texttt{pynescript-lsp}) with multi-editor packaging. The design goal is offline IDE
parity for a financial \textsc{Dsl}: full local language services without mandatory
network calls, while cleanly separating optional evaluation and chart/backtest
surfaces (Pro \api{}) that share the same \ast{} and runtime.

\paragraph{Contributions.}
This paper makes the following contributions:
\begin{enumerate}[leftmargin=*]
  \item A complete \lsp{} architecture for a bar-oriented financial \textsc{Dsl},
        including workspace document state, cached parse/lint, and feature modules
        that take pure \((params, source)\) inputs rather than holding server globals
        (Section~\ref{sec:architecture}).
  \item Detailed designs for diagnostics, completion over \(\sim\)800+ builtin metadata
        items, hover, navigation, semantic tokens, inlay hints, and unparse-based
        formatting (Section~\ref{sec:features}).
  \item A metadata pipeline with generation, optional Fernet encryption, SHA-256
        integrity, and dual-path runtime load for development versus Nuitka
        distributions (Section~\ref{sec:metadata}).
  \item Multi-editor delivery: a TypeScript VS~Code extension with TextMate grammar
        and bundled language client, plus thin configs for Neovim, Zed, and Emacs
        (Section~\ref{sec:editors}).
  \item A complementary \cli{} surface and a clear boundary with evaluation / Pro
        \api{} endpoints (Sections~\ref{sec:cli}--\ref{sec:eval}).
  \item Implementation lessons, limitations, and related work for language tooling
        over financial \textsc{Dsl}s (Sections~\ref{sec:lessons}--\ref{sec:conclusion}).
\end{enumerate}

% =============================================================================
\section{Background}
\label{sec:background}
% =============================================================================

\subsection{Language Server Protocol}
\label{sec:bg-lsp}

The \lsp{}~\citep{lsp2023} standardizes a set of methods over JSON-RPC, typically
transported on STDIO between an editor process and a long-lived server. Core concepts
include:
\begin{itemize}[leftmargin=*]
  \item \textbf{Lifecycle:} \texttt{initialize} / \texttt{initialized} / \texttt{shutdown}
        / \texttt{exit}, with capability negotiation.
  \item \textbf{Document synchronization:} \texttt{textDocument/didOpen},
        \texttt{didChange} (full or incremental), \texttt{didSave}, \texttt{didClose}.
  \item \textbf{Language features:} completion, hover, definition, references, symbols,
        formatting, diagnostics (push via \texttt{publishDiagnostics} and pull via
        \texttt{textDocument/diagnostic} in 3.16+), semantic tokens, and inlay hints.
\end{itemize}
By moving language intelligence into a server, the \(M \times N\) problem of \(M\)
languages and \(N\) editors collapses to \(M + N\)~\citep{voutilainen2016lsp}. For
\textsc{Dsl}s with modest user bases, this amortization is essential: one implementation
can serve many editor communities.

\subsection{pygls and the Python LSP stack}
\label{sec:bg-pygls}

\texttt{pygls} (Python Generic Language Server)~\citep{murphy2020pygls} provides a
decorator-oriented API for registering \lsp{} method handlers on a
\texttt{LanguageServer} subclass, with typed messages via \texttt{lsprotocol}. It is a
natural fit for \pyne{} because the entire language pipeline (parser, \ast{}, linter,
unparser, evaluator) is already Python. Alternatives such as writing a server in
TypeScript or Rust would require re-implementing or FFI-bridging the grammar and
analysis stack. Compiling the server to a Nuitka onefile binary~\citep{nuitka} further
reduces the end-user dependency surface for editor extensions.

\subsection{Pine Script editing pain points}
\label{sec:bg-pain}

Empirical pain points for offline \pine{} development include:
\begin{enumerate}[leftmargin=*]
  \item \textbf{No offline diagnostics.} Syntax and many semantic mistakes surface only
        after paste into a hosted editor or after a failed chart compile.
  \item \textbf{Builtin discoverability.} Hundreds of functions across namespaces
        (\texttt{ta.*}, \texttt{strategy.*}, \texttt{request.*}, drawing objects, \ldots)
        are hard to recall without hover and completion.
  \item \textbf{Formatting drift.} Teams without a shared pretty-printer accumulate
        style noise in Git diffs.
  \item \textbf{Editor lock-in.} Browser-only editing precludes Neovim/Emacs power users
        and headless CI checks.
  \item \textbf{Separation of editing and evaluation.} Chart preview and backtest are
        valuable but should not be prerequisites for static language services.
\end{enumerate}
\pyne{} addresses (1)--(4) with \texttt{pynescript-lsp} and the \cli{}; (5) is
intentionally a separate process boundary (Section~\ref{sec:eval}).

\subsection{Language tooling for DSLs}
\label{sec:bg-dsl}

Language workbenches and embedded \textsc{Dsl}s have long argued for first-class
tooling~\citep{fowler2010dsl,hudak1996dsl}. Interpreter textbooks emphasize the
centrality of a well-defined \ast{} for analysis and transformation~\citep{nystrom2021crafting}.
\pyne{} follows this orthodoxy: the \lsp{} is a consumer of the same \ast{} used by
evaluation, not a parallel frontend.

% =============================================================================
\section{System Architecture of \texttt{pynescript-lsp}}
\label{sec:architecture}
% =============================================================================

\subsection{High-level topology}
\label{sec:arch-topo}

Figure~\ref{fig:topology} summarizes the multi-surface deployment. Editors speak \lsp{}
over STDIO to \texttt{pynescript-lsp}. The server owns a \texttt{Workspace} of open
documents, re-parses and re-lints on mutation, and dispatches feature handlers. The open
package (\texttt{src/pynescript/}) supplies parser, \ast{}, linter, and unparser. A
Pro \api{} (Flask, optionally Cloud Run) shares the evaluator for chart preview and
backtest but is not on the \lsp{} critical path.

\begin{figure}[t]
\centering
\begin{tikzpicture}[
  font=\small,
  box/.style={draw, rounded corners=2pt, align=center, inner sep=6pt, minimum height=1.1cm},
  ed/.style={box, fill=pyneblue!12, minimum width=2.6cm},
  srv/.style={box, fill=pyneorange!15, minimum width=3.4cm},
  core/.style={box, fill=pynegreen!12, minimum width=2.8cm},
  pro/.style={box, fill=pynegray!18, minimum width=2.6cm},
  arr/.style={-{Latex}, thick}
]
  \node[ed] (vscode) {VS~Code\\extension};
  \node[ed, right=0.45cm of vscode] (nvim) {Neovim};
  \node[ed, right=0.45cm of nvim] (zed) {Zed};
  \node[ed, right=0.45cm of zed] (emacs) {Emacs};
  \node[srv, below=1.1cm of nvim, xshift=0.9cm] (lsp) {\texttt{pynescript-lsp}\\(pygls + Nuitka)};
  \node[core, below left=1.2cm and 0.3cm of lsp] (ast) {Parser / AST\\Linter / Unparser};
  \node[core, below right=1.2cm and 0.3cm of lsp] (meta) {Builtin metadata\\Completion providers};
  \node[pro, right=1.6cm of lsp] (api) {Pro API\\(eval / chart / backtest)};

  \draw[arr] (vscode) -- (lsp);
  \draw[arr] (nvim) -- (lsp);
  \draw[arr] (zed) -- (lsp);
  \draw[arr] (emacs) -- (lsp);
  \draw[arr] (lsp) -- (ast);
  \draw[arr] (lsp) -- (meta);
  \draw[arr, dashed] (api) -- node[above, font=\scriptsize, sloped] {shared \ast{}/runtime} (ast);
\end{tikzpicture}
\caption{Multi-editor deployment topology. All editors use STDIO \lsp{}; the Pro \api{}
shares \ast{}/runtime but is a separate process and auth model.}
\label{fig:topology}
\end{figure}

\subsection{Server process model}
\label{sec:arch-server}

\texttt{PynescriptLanguageServer} subclasses \texttt{pygls.lsp.server.LanguageServer}
with name \texttt{"Pynescript"} and the package version. On construction it creates an
empty \texttt{Workspace} and registers handlers in \texttt{setup\_method\_handlers()}.
Entry points are:
\begin{itemize}[leftmargin=*]
  \item Console script \texttt{pynescript-lsp} \(\rightarrow\)
        \texttt{pynescript.langserver.\_\_main\_\_:main}
  \item Module form \texttt{python -m pynescript.langserver}
  \item Optional Nuitka binary under \texttt{dist/lsp/pynescript-lsp}
\end{itemize}
Transport defaults to STDIO (\texttt{start\_io()}), matching editor integrations.

\subsection{Workspace and document state}
\label{sec:arch-workspace}

Each open URI maps to a \texttt{TextDocumentState}:
\begin{center}
\small
\begin{tabular}{ll}
\toprule
Field & Role \\
\midrule
\texttt{uri}, \texttt{source}, \texttt{version} & Buffer identity and content \\
\texttt{ast} & Cached parse tree, or \texttt{None} on failure \\
\texttt{diagnostics} & List of \texttt{LintWarning} (not yet \lsp{} types) \\
\texttt{parse\_error}, \texttt{parse\_error\_line} & E001 synthesis inputs \\
\bottomrule
\end{tabular}
\end{center}
On \texttt{didOpen}/\texttt{didChange}, the workspace applies full-document or
incremental edits, then runs \texttt{\_parse\_and\_lint}: successful parse feeds
\texttt{lint\_script}; exceptions clear the \ast{} and store a human-readable error.
Feature handlers typically read \texttt{get\_source(uri)} and optionally receive a
pre-parsed tree to avoid redundant work. The buffer---not disk---is the source of
truth until the client saves.

\subsection{Capability declaration}
\label{sec:arch-caps}

Capabilities are declared once from \texttt{config.get\_server\_capabilities()} during
\texttt{initialize}. Table~\ref{tab:methods} lists supported methods. Signature help and
code actions are intentionally omitted until first-class handlers exist; advertising
unimplemented providers would mislead clients.

\begin{table}[t]
\centering
\caption{\lsp{} methods implemented by \texttt{pynescript-lsp}.}
\label{tab:methods}
\small
\begin{tabular}{@{}lll@{}}
\toprule
Method & Capability notes & Module \\
\midrule
\texttt{initialize} / \texttt{initialized} / \texttt{shutdown} & Lifecycle & \texttt{server.py} \\
\texttt{textDocument/didOpen|Change|Close|Save} & Incremental sync & \texttt{server.py} \\
\texttt{textDocument/publishDiagnostics} & Push on open/change/save & \texttt{workspace} \\
\texttt{textDocument/diagnostic} & Pull (full report + \texttt{result\_id}) & \texttt{server.py} \\
\texttt{workspace/diagnostic} & Per open document & \texttt{server.py} \\
\texttt{textDocument/completion} & Trigger \texttt{"."}; resolve provider & \texttt{completion.py} \\
\texttt{completionItem/resolve} & Docs rehydration & \texttt{completion.py} \\
\texttt{textDocument/hover} & Markdown signatures + docs & \texttt{hover.py} \\
\texttt{textDocument/definition} & Same-document & \texttt{definitions.py} \\
\texttt{textDocument/references} & Same-document & \texttt{references.py} \\
\texttt{textDocument/documentSymbol} & Outline & \texttt{symbols.py} \\
\texttt{workspace/symbol} & Query over open docs & \texttt{symbols.py} \\
\texttt{textDocument/formatting} & Full document & \texttt{formatting.py} \\
\texttt{textDocument/rangeFormatting} & Range unparse & \texttt{formatting.py} \\
\texttt{textDocument/semanticTokens/full} & Delta-encoded legend & \texttt{semantic\_tokens.py} \\
\texttt{textDocument/inlayHint} & Inferred declaration types & \texttt{inlay\_hints.py} \\
\bottomrule
\end{tabular}
\end{table}

\subsection{Request lifecycle}
\label{sec:arch-lifecycle}

Figure~\ref{fig:lifecycle} shows a representative completion request after document
open. Diagnostics are published eagerly; feature methods are pull-style request/response.

\begin{figure}[t]
\centering
\begin{tikzpicture}[
  font=\scriptsize,
  actor/.style={font=\small\bfseries},
  msg/.style={-{Latex}, thick},
  rmsg/.style={-{Latex}, thick, dashed}
]
  \node[actor] (C) at (0,0) {Editor client};
  \node[actor] (S) at (5.2,0) {Language server};
  \node[actor] (W) at (10.2,0) {Workspace};

  \draw[thick] (0,-0.3) -- (0,-7.6);
  \draw[thick] (5.2,-0.3) -- (5.2,-7.6);
  \draw[thick] (10.2,-0.3) -- (10.2,-7.6);

  \draw[msg] (0,-0.8) -- node[above] {\texttt{initialize}} (5.2,-0.8);
  \draw[rmsg] (5.2,-1.3) -- node[above] {\texttt{InitializeResult} + caps} (0,-1.3);
  \draw[msg] (0,-1.9) -- node[above] {\texttt{initialized}} (5.2,-1.9);

  \draw[msg] (0,-2.6) -- node[above] {\texttt{didOpen}(uri, text)} (5.2,-2.6);
  \draw[msg] (5.2,-3.1) -- node[above] {\texttt{put\_document}} (10.2,-3.1);
  \draw[rmsg] (10.2,-3.6) -- node[above] {parse + lint} (5.2,-3.6);
  \draw[rmsg] (5.2,-4.2) -- node[above] {\texttt{publishDiagnostics}} (0,-4.2);

  \draw[msg] (0,-5.0) -- node[above] {\texttt{completion}(uri, pos)} (5.2,-5.0);
  \draw[msg] (5.2,-5.6) -- node[above] {\texttt{get\_source}} (10.2,-5.6);
  \draw[rmsg] (10.2,-6.2) -- node[above] {source text} (5.2,-6.2);
  \draw[rmsg] (5.2,-6.9) -- node[above] {\texttt{CompletionList}} (0,-6.9);
\end{tikzpicture}
\caption{\lsp{} request lifecycle: initialize, document open with push diagnostics, then
a completion request served from workspace source and metadata providers.}
\label{fig:lifecycle}
\end{figure}

\subsection{Feature module layout}
\label{sec:arch-modules}

Figure~\ref{fig:modules} depicts the package structure under
\texttt{src/pynescript/langserver/}. Handlers in \texttt{features/} are pure-ish
functions; \texttt{providers/} own metadata load and \texttt{CompletionItem} construction;
\texttt{protocol/} holds word/trigger utilities; \texttt{config.py} is the single
capability table.

\begin{figure}[t]
\centering
\begin{tikzpicture}[
  font=\small,
  node distance=0.55cm and 0.7cm,
  box/.style={draw, rounded corners=2pt, align=center, inner sep=5pt, fill=white},
  root/.style={box, fill=pyneorange!20, font=\small\bfseries},
  mid/.style={box, fill=pyneblue!12},
  leaf/.style={box, fill=codebg, font=\scriptsize}
]
  \node[root] (root) {\texttt{langserver/}};
  \node[mid, below left=1.0cm and 2.4cm of root] (feat) {\texttt{features/}};
  \node[mid, below=1.0cm of root] (prov) {\texttt{providers/}};
  \node[mid, below right=1.0cm and 2.4cm of root] (proto) {\texttt{protocol/}};
  \node[mid, right=0.5cm of prov] (ws) {\texttt{workspace.py}};
  \node[mid, left=0.5cm of prov] (srv) {\texttt{server.py}};
  \node[mid, above right=0.15cm and 0.1cm of srv] (cfg) {\texttt{config.py}};

  \node[leaf, below=0.45cm of feat] (f1) {completion, diagnostics, hover\\formatting, definitions, references\\semantic\_tokens, symbols, inlay\_hints};
  \node[leaf, below=0.45cm of prov] (p1) {builtin\_metadata(.json/.enc)\\completion\_items, metadata\_decrypt};
  \node[leaf, below=0.45cm of proto] (pr1) {utils, constants};

  \draw[-{Latex}] (root) -- (feat);
  \draw[-{Latex}] (root) -- (prov);
  \draw[-{Latex}] (root) -- (proto);
  \draw[-{Latex}] (root) -- (srv);
  \draw[-{Latex}] (root) -- (ws);
  \draw[-{Latex}] (root) -- (cfg);
  \draw[-{Latex}] (feat) -- (f1);
  \draw[-{Latex}] (prov) -- (p1);
  \draw[-{Latex}] (proto) -- (pr1);
\end{tikzpicture}
\caption{Feature module architecture of \texttt{pynescript-lsp}.}
\label{fig:modules}
\end{figure}

% =============================================================================
\section{Feature Design}
\label{sec:features}
% =============================================================================

\subsection{Diagnostics pipeline}
\label{sec:feat-diag}

Diagnostics are the primary static feedback channel. Figure~\ref{fig:diagnostics}
shows the pipeline. On every open, change, and save the workspace re-parses. Parse
success feeds the orthogonal static linter (\texttt{lint\_script}); parse failure
synthesizes diagnostic code \texttt{E001} at the extracted line. Findings convert to
\texttt{lsp.Diagnostic} with \texttt{source="PineScript"}, severity mapping
(error/warning/info/hint), and optional tags (e.g., unnecessary/deprecated).

\begin{figure}[t]
\centering
\begin{tikzpicture}[
  font=\small,
  node distance=0.85cm,
  flowstep/.style={draw, rounded corners=2pt, align=center, inner sep=6pt, minimum width=3.2cm, fill=codebg},
  dec/.style={draw, diamond, aspect=2.2, align=center, inner sep=1pt, fill=pyneorange!12},
  arr/.style={-{Latex}, thick}
]
  \node[flowstep] (src) {Buffer source};
  \node[flowstep, below=of src] (parse) {\texttt{parse(source)}};
  \node[dec, below=of parse] (ok) {OK?};
  \node[flowstep, below left=1.1cm and 0.6cm of ok] (e001) {E001 parse error};
  \node[flowstep, below right=1.1cm and 0.6cm of ok] (lint) {\texttt{lint\_script}};
  \node[flowstep, below=2.6cm of ok] (map) {Map to \lsp{} diagnostics};
  \node[flowstep, below=of map] (pub) {\texttt{publishDiagnostics} / pull};

  \draw[arr] (src) -- (parse);
  \draw[arr] (parse) -- (ok);
  \draw[arr] (ok) -- node[left, font=\scriptsize] {no} (e001);
  \draw[arr] (ok) -- node[right, font=\scriptsize] {yes} (lint);
  \draw[arr] (e001) |- (map);
  \draw[arr] (lint) |- (map);
  \draw[arr] (map) -- (pub);
\end{tikzpicture}
\caption{Diagnostics pipeline: parse, optional lint, conversion, push/pull publish.}
\label{fig:diagnostics}
\end{figure}

Representative lint codes include \texttt{E001} (parse), \texttt{W001}/\texttt{W002}
(warnings), and \texttt{C001}--\texttt{C004} (style: complexity/length/if-style/newline).
Push diagnostics use \texttt{textDocument/publishDiagnostics}; pull clients may call
\texttt{textDocument/diagnostic}, receiving a full report with
\texttt{result\_id=}\textit{uri}\texttt{-}\textit{version}. Algorithm~\ref{alg:publish} formalizes
the push path.

\begin{algorithm}[t]
\caption{PublishDiagnostics on document mutation}
\label{alg:publish}
\begin{algorithmic}[1]
\Require URI \(u\), source \(s\), version \(v\)
\State \(doc \gets \Call{PutOrUpdate}{u, s, v}\)
\State \Call{ParseAndLint}{doc}
\State \(D \gets [\,]\)
\If{\(doc.\mathit{parse\_error} \neq \bot\)}
  \State append E001 diagnostic at \(doc.\mathit{parse\_error\_line}\) to \(D\)
\EndIf
\For{each warning \(w\) in \(doc.\mathit{diagnostics}\)}
  \State append \Call{ToLspDiagnostic}{w, s} to \(D\)
\EndFor
\State \Call{PublishDiagnostics}{u, D}
\end{algorithmic}
\end{algorithm}

\subsection{Completion from builtin metadata}
\label{sec:feat-completion}

Completion is metadata-driven, not evaluator-driven. The server loads a catalog of
builtins (plaintext JSON in development, Fernet-encrypted in compiled binaries;
Section~\ref{sec:metadata}), filters by prefix or module path, and returns
\texttt{CompletionItem} rows. Current catalog size is \(\sim\)872 entries across 25
categories (e.g., \texttt{ta.technical\_analysis}, \texttt{strategy}, \texttt{array},
\texttt{request}, drawing objects). Historically documentation referred to
\(\sim\)482 builtins; the completion surface is larger because module members, aliases,
and related symbols expand the item set. We report ``hundreds of builtins, \(\sim\)800+
completion items/metadata.''

Figure~\ref{fig:completion} illustrates data flow from metadata artifacts through
providers to the client.

\begin{figure}[t]
\centering
\begin{tikzpicture}[
  font=\small,
  node distance=0.7cm and 0.9cm,
  box/.style={draw, rounded corners=2pt, align=center, inner sep=5pt, fill=codebg},
  arr/.style={-{Latex}, thick}
]
  \node[box, fill=pynegreen!12] (json) {\texttt{builtin\_metadata.json}\\(+ .enc / .sha256)};
  \node[box, right=of json] (load) {\texttt{get\_metadata()}\\cache};
  \node[box, right=of load] (ctx) {Cursor context\\prefix / trigger \texttt{.}};
  \node[box, below=1.1cm of load] (filt) {\texttt{fuzzy\_filter} /\\module filter};
  \node[box, below=1.1cm of ctx] (build) {\texttt{build\_completion\_*}\\snippets, sort\_text};
  \node[box, below=1.1cm of filt, fill=pyneblue!12] (list) {\texttt{CompletionList}};
  \node[box, right=of list, fill=pyneorange!15] (client) {Editor UI};

  \draw[arr] (json) -- (load);
  \draw[arr] (load) -- (ctx);
  \draw[arr] (load) -- (filt);
  \draw[arr] (ctx) -- (build);
  \draw[arr] (filt) -- (build);
  \draw[arr] (build) -- (list);
  \draw[arr] (list) -- (client);
\end{tikzpicture}
\caption{Completion data flow: metadata \(\rightarrow\) items \(\rightarrow\) client.}
\label{fig:completion}
\end{figure}

\paragraph{Context and ranking.}
Algorithm~\ref{alg:completion} sketches request handling. Trigger character \texttt{.}
and dotted prefixes (\texttt{ta.sm}) route to module completion; otherwise a fuzzy filter
scores candidates (exact label 1000, prefix 500, substring 100, category 50, brief 25;
default limit 50). Sort keys prefer modules (\(\texttt{\textbackslash x01}\ldots\)) over
root symbols (\(\texttt{\textbackslash x02}\ldots\)). Category headers appear as
folder-kind items with empty insert text. Snippets use \texttt{InsertTextFormat.Snippet}
when metadata supplies \texttt{\$\{...\}} placeholders. Resolve rehydrates full
documentation for a single item via \texttt{get\_builtin(label)}.

\begin{algorithm}[t]
\caption{HandleCompletion}
\label{alg:completion}
\begin{algorithmic}[1]
\Require params \(P\), source \(s\)
\State \(t \gets\) text before cursor on line of \(P.\mathit{position}\)
\State \(prefix \gets\) last token of \(t\)
\If{\texttt{"."} \(\in prefix\) \textbf{or} trigger char is \texttt{"."}}
  \State \(m \gets\) module name before final \texttt{.}
  \State \Return \Call{BuildModuleCompletion}{m}
\Else
  \State \Return \Call{BuildCompletionList}{prefix}
\EndIf
\end{algorithmic}
\end{algorithm}

\subsection{Hover and signatures}
\label{sec:feat-hover}

Hover resolves the word under the cursor---including dotted forms such as
\texttt{ta.sma}---against \texttt{get\_builtin}. Successful hits return Markdown
combining signature \texttt{detail}, brief, and extended documentation. The server does
not yet advertise a dedicated \texttt{signatureHelp} provider; signature text is instead
surfaced through hover and completion \texttt{detail} fields, which is adequate for
most Pine call sites and avoids half-implemented capability flags.

\subsection{Navigation: definitions, references, symbols}
\label{sec:feat-nav}

\textbf{Definition} walks the \ast{} with a \texttt{DefinitionFinder} visitor for the
identifier under the cursor and returns same-document \texttt{Location}s (functions,
types, assignments). \textbf{References} symmetrically collect name occurrences.
\textbf{Document symbols} build an outline (functions \(\rightarrow\) Function, type
defs \(\rightarrow\) Class, assignments \(\rightarrow\) Variable). \textbf{Workspace
symbols} scan all open documents and filter by case-insensitive substring query.
Cross-file project indexes and library \texttt{import} resolution are future work
(Section~\ref{sec:limits}).

\subsection{Semantic tokens and inlay hints}
\label{sec:feat-sem}

Semantic tokens extend TextMate grammars with \ast{}-backed classification. The server
advertises a legend of thirteen types (namespace, type, class, function, method,
variable, parameter, property, keyword, string, number, operator, comment) and four
modifiers (declaration, definition, readonly, defaultLibrary). Tokens are delta-encoded
for \texttt{textDocument/semanticTokens/full}; range requests are not advertised.
Builtin namespaces (\texttt{ta}, \texttt{math}, \texttt{strategy}, \ldots) receive
\texttt{defaultLibrary} modifiers.

Inlay hints show confident, shape-based inferred types next to declarations, e.g.
\texttt{length = 14} \(\rightarrow\) \texttt{: const int}, \texttt{rsi = ta.rsi(...)}
\(\rightarrow\) \texttt{: series float}, drawing return types from metadata
\texttt{detail} after \(\rightarrow\). Only high-confidence cases are emitted
(\texttt{resolve\_provider=False}).

\subsection{Formatting via unparse}
\label{sec:feat-format}

Formatting reuses the lossless unparser rather than a separate pretty-print grammar
(Algorithm~\ref{alg:format}). Full-document and range formatting both parse, unparse,
and emit a single replace \texttt{TextEdit} when the result differs. Parse failures
return an empty edit list (no throw to the client), so broken mid-edit buffers never
hard-fail format-on-save. This design keeps formatting semantics aligned with the
canonical \ast{} round-trip used in tests and \cli{} \texttt{format}.

\begin{algorithm}[t]
\caption{FormatDocument}
\label{alg:format}
\begin{algorithmic}[1]
\Require source \(s\)
\If{\(s = \bot\) or \(s = \varepsilon\)}
  \State \Return \([\,]\)
\EndIf
\State \textbf{try} \(tree \gets \Call{Parse}{s}\)
\State \(s' \gets \Call{Unparse}{tree}\)
\If{\(s' = s\)}
  \State \Return \([\,]\)
\EndIf
\State \Return \([\Call{ReplaceEntireDocument}{s'}]\)
\State \textbf{catch} \Return \([\,]\)
\end{algorithmic}
\end{algorithm}

% =============================================================================
\section{Metadata Pipeline}
\label{sec:metadata}
% =============================================================================

\subsection{Generation}
\label{sec:meta-gen}

Builtin metadata is a map from fully qualified names (e.g., \texttt{ta.sma},
\texttt{strategy.entry}) to records:
\begin{lstlisting}[style=python,numbers=none,basicstyle=\ttfamily\scriptsize]
{
  "ta.sma": {
    "label": "ta.sma",
    "kind": "function",
    "detail": "ta.sma(series, int) -> series float",
    "brief": "...",
    "documentation": "...",
    "snippet": "ta.sma(${1:param1}, ${2:param2})",
    "category": "ta.technical_analysis"
  }
}
\end{lstlisting}
The generator \texttt{scripts/generate\_builtin\_metadata.py} introspects evaluator
builtins and related surfaces; the JSON under
\texttt{langserver/providers/} is regenerated rather than hand-edited as source of
truth. Consumers include completion, hover, and inlay return-type extraction.

\subsection{Encryption and integrity}
\label{sec:meta-enc}

For Nuitka-shipped binaries, plaintext is not ideal: casual
\texttt{strings pynescript-lsp} would dump the documentation corpus. The build
optionally encrypts with Fernet (AES-128-CBC + HMAC-SHA256) to
\texttt{builtin\_metadata.json.enc}, accompanied by a 16-character SHA-256 prefix of
the plaintext in \texttt{builtin\_metadata.json.sha256}. The key is supplied via
\texttt{CRYPTO\_KEY} / \texttt{.metadata.key} in CI and is not committed. This is
\emph{distribution hygiene / obscurity}, not an access-control boundary: anyone with
the binary can eventually recover the catalog. It preserves a future knob for paid
metadata endpoints without changing handler code.

\subsection{Runtime load order}
\label{sec:meta-load}

\texttt{get\_metadata()} uses a process-local cache and the following order:
\begin{enumerate}[leftmargin=*]
  \item Prefer plaintext \texttt{builtin\_metadata.json} when present and valid
        (development always wins after regen).
  \item Else decrypt \texttt{.enc} via \texttt{metadata\_decrypt.load\_encrypted\_metadata()},
        validating the SHA-256 prefix when present.
  \item Else return an empty dict (logged warnings); completion degrades gracefully.
\end{enumerate}
Key resolution checks \texttt{.metadata.key} beside the package (or under
\texttt{sys.\_MEIPASS} when frozen), then \texttt{PYNESCRIPT\_METADATA\_KEY}.

% =============================================================================
\section{Multi-Editor Delivery}
\label{sec:editors}
% =============================================================================

\subsection{VS Code extension}
\label{sec:vscode}

The \texttt{vscode-extension/} package (\textbf{PYNE}, publisher
\texttt{jango-blockchained}) is a TypeScript language client that:
\begin{itemize}[leftmargin=*]
  \item Registers language id \texttt{pinescript} with extensions \texttt{.pyne},
        \texttt{.pine}, \texttt{.pinev5}, \texttt{.pinev6}, \texttt{.pinescript}.
  \item Contributes TextMate grammar \texttt{syntaxes/pinescript.tmLanguage.json}
        for baseline highlighting.
  \item Auto-detects the server: \texttt{pynescript-lsp} on \texttt{PATH}, else
        \texttt{python -m pynescript.langserver}, else a configured absolute path.
  \item Enables semantic highlighting defaults, sets the extension as default
        formatter, and exposes settings for diagnostics/completion/snippets.
\end{itemize}
Install path for users: \texttt{pip install "hoox-pyne[lsp]"} plus the VSIX, or a
future marketplace listing. The extension does not reimplement language intelligence;
it is a thin client over STDIO.

\subsection{Thin clients: Neovim, Zed, Emacs}
\label{sec:thin}

Directory \texttt{clients/} ships ready snippets:
\begin{itemize}[leftmargin=*]
  \item \textbf{Neovim} (\texttt{neovim.lua}): \texttt{nvim-lspconfig} registration,
        keymaps for definition, references, hover, format.
  \item \textbf{Zed} (\texttt{zed.json}): language server command
        \texttt{["pynescript-lsp"]} with STDIO.
  \item \textbf{Emacs} (\texttt{emacs.el}): \texttt{lsp-mode} client registration for
        \texttt{pinescript-mode}.
\end{itemize}
Helix and Sublime Text configurations are documented similarly. Any STDIO-capable
client can launch \texttt{pynescript-lsp --stdio}. Table~\ref{tab:editors} summarizes
the matrix.

\begin{table}[t]
\centering
\caption{Editor delivery matrix for \pyne{} language services.}
\label{tab:editors}
\small
\begin{tabular}{@{}llll@{}}
\toprule
Editor & Client form & Server discovery & Syntax \\
\midrule
VS~Code / Cursor & TypeScript extension (VSIX) & auto / PATH / module & TextMate + semantic \\
Neovim & \texttt{clients/neovim.lua} & PATH binary or module & Treesitter optional \\
Zed & \texttt{clients/zed.json} & PATH & Zed grammar optional \\
Emacs & \texttt{clients/emacs.el} & PATH & major mode \\
Helix / Sublime & docs snippets & PATH & client-specific \\
\bottomrule
\end{tabular}
\end{table}

\subsection{Packaging channels}
\label{sec:packaging}

Three delivery shapes share one codebase~\citep{pynescript2026}:
\begin{enumerate}[leftmargin=*]
  \item \textbf{pip:} \texttt{hoox-pyne[lsp]} for libraries, \cli{}, and module launch.
  \item \textbf{Nuitka onefile:} self-contained \texttt{pynescript-lsp} for editors that
        prefer a single executable~\citep{nuitka}.
  \item \textbf{VS~Code VSIX:} Node client + server discovery; optional binary bundle.
\end{enumerate}
Container images on GHCR cover \cli{}/\api{} operational use cases and are orthogonal
to interactive \lsp{} sessions.

% =============================================================================
\section{CLI as Complementary Tooling Surface}
\label{sec:cli}
% =============================================================================

The console script \texttt{pynescript} (Click-based) is deliberately \emph{not} a
subcommand of the language server. It targets headless and CI workflows that
complement the editor:

\begin{center}
\small
\begin{tabular}{@{}ll@{}}
\toprule
Command & Role \\
\midrule
\texttt{check} & Parse-only validation (CI-friendly exit codes) \\
\texttt{format} / \texttt{fmt} & Parse \(\rightarrow\) unparse; optional in-place write \\
\texttt{lint} & Static rules; colored or JSON output \\
\texttt{parse-and-dump} / \texttt{dump} / \texttt{ast} & AST dump \\
\texttt{parse-and-unparse} / \texttt{unparse} & Round-trip source \\
\texttt{compile} & Transpile / Numba compile-check \\
\texttt{run} & Compile + execute on synthetic \ohlcv{} \\
\texttt{data} & Fetch market bars (mock / Yahoo / \ldots) \\
\texttt{prewarm} & Warm Numba / IR caches \\
\texttt{info} & Version and optional extras \\
\bottomrule
\end{tabular}
\end{center}

Shared components ensure behavioral parity: \cli{} \texttt{format} and \lsp{} formatting
both call the unparser; \cli{} \texttt{lint} and \lsp{} diagnostics both call
\texttt{lint\_script}. The language server remains a separate entry point
(\texttt{pynescript-lsp}) so editor sessions do not inherit \cli{} flag parsing or
network data providers.

% =============================================================================
\section{Integration with Evaluation}
\label{sec:eval}
% =============================================================================

\subsection{Shared AST, separate processes}
\label{sec:eval-boundary}

Evaluation (literal expressions, full scripts, strategy backtests) uses the same
parse \(\rightarrow\) \ast{} path as the language server but is \emph{not} invoked on
every keystroke. The design document~\citep{pynescript2026} makes this separation
explicit: lint is a cheap static phase; evaluation is a different cost profile and
failure surface.

\subsection{Pro API surface}
\label{sec:eval-pro}

The Pro \api{} (\texttt{backend/}, Flask) exposes HTTP endpoints such as evaluate, chart
preview, and quick backtest. Auth (API keys, tiers, rate limits) and scaling (gunicorn /
Cloud Run) differ completely from single-user STDIO. Taking the \api{} down must not
break offline editing. Optionally, an HTTP bridge reuses completion/hover handlers for
browser-hosted tools (AXIS) that cannot spawn pygls in-tab; that bridge is a convenience
facade, not a replacement for the full \lsp{}.

\subsection{IDE-perspective run/preview}
\label{sec:eval-ide}

From an IDE user's perspective, optional ``run script'' or ``preview chart'' actions
are natural extensions. The preferred architecture is: editor command \(\rightarrow\)
local \cli{} or authenticated \api{} call \(\rightarrow\) render result in a webview or
output channel---never by embedding market I/O inside diagnostic hot paths. This keeps
the open-source editor experience fully offline while allowing premium evaluation
tiers~\citep{hoox2026}.

% =============================================================================
\section{Implementation Experience and Design Lessons}
\label{sec:lessons}
% =============================================================================

\paragraph{Advertise only what you implement.}
Early drafts considered declaring signature help and code actions ``for free.'' Clients
then surface broken UI affordances. \pyne{} now lists only wired handlers in
\texttt{get\_server\_capabilities()}.

\paragraph{Pure feature handlers beat god-objects.}
Handlers of the form \texttt{handle\_X(params, source)} (and optional cached tree) are
unit-testable without spinning STDIO. Server methods become thin adapters.

\paragraph{Parse failure must be non-fatal.}
Mid-edit buffers are often temporarily invalid. The server retains process health,
returns empty navigation/completion when needed, and still publishes E001.

\paragraph{Unparse is a formatter.}
For \textsc{Dsl}s with a maintained unparser, a separate formatting grammar is
unnecessary duplication. Round-trip tests become formatting regression tests.

\paragraph{Metadata is a product artifact.}
Treating builtin docs as generated, versioned, optionally encrypted data---not as
hard-coded strings in handlers---keeps completion/hover in sync with the evaluator
surface and enables binary packaging policies.

\paragraph{Multi-editor thin clients scale.}
Investing in a solid STDIO server plus a polished VS~Code client, while shipping
minimal configs for other editors, matches real adoption curves without \(N\) full
extensions.

\paragraph{Financial DSL quirks.}
Pine's series semantics, \texttt{var}/\texttt{varip}, and version annotations affect
static analysis more than completion. Inlay hints remain intentionally conservative
(shape-based) rather than claiming full type inference parity with the hosted
compiler~\citep{nystrom2021crafting,pierce2002types}.

% =============================================================================
\section{Related Work}
\label{sec:related}
% =============================================================================

\paragraph{LSP ecosystem.}
The \lsp{} specification~\citep{lsp2023} and early advocacy~\citep{voutilainen2016lsp}
established the multi-editor language service pattern. Production servers for
TypeScript, Rust, Go, and Python demonstrate capability breadth; \texttt{pygls}
lowers the barrier for Python-implemented servers~\citep{murphy2020pygls}.
\texttt{pynescript-lsp} applies this stack to a financial \textsc{Dsl} rather than a
general-purpose language.

\paragraph{Language workbenches and DSLs.}
Fowler~\citep{fowler2010dsl} and Hudak~\citep{hudak1996dsl} emphasize that
\textsc{Dsl} value depends on tooling as much as semantics. Workbenches often generate
editors from language definitions; \pyne{} instead reuses a hand-maintained ANTLR4
grammar~\citep{parr2013antlr,parr2011allstar} and ASDL schema~\citep{wang1997asdl}
shared with the runtime---closer to interpreter engineering~\citep{nystrom2021crafting}
than to generative workbenches.

\paragraph{Financial IDE tooling.}
Hosted trading platforms provide tightly integrated script editors with chart
feedback but limited offline/CI support. Quantitative libraries (e.g., Python
data stacks~\citep{brooks2015python,harris2020numpy}) offer powerful analysis without
Pine-specific language services. \pyne{} targets the gap: local language intelligence
for \pine{} sources plus optional evaluation, without claiming affiliation with the
platform vendor~\citep{tradingview-pine}.

\paragraph{Binary distribution of Python tools.}
Nuitka~\citep{nuitka} and similar compilers address cold-start and dependency issues
for shipping Python language servers to users who should not manage virtualenvs. Our
encrypted metadata path is a packaging concern orthogonal to language semantics.

% =============================================================================
\section{Limitations and Future Work}
\label{sec:limits}
% =============================================================================

\begin{itemize}[leftmargin=*]
  \item \textbf{Same-document navigation.} Definitions/references do not yet resolve
        cross-file \texttt{import} graphs or library version pins.
  \item \textbf{No signature help / code actions.} Documented omissions; quick-fix
        stubs exist in diagnostics helpers but are not first-class features.
  \item \textbf{Incremental parse.} The workspace re-parses whole buffers; tree-sitter
        or incremental ANTLR strategies could reduce latency on large scripts.
  \item \textbf{Type system depth.} Inlay hints are shape-based; full series/type
        inference and strict v6 boolean semantics in the language server remain
        incomplete relative to the evaluator.
  \item \textbf{Hosted editor parity.} Chart-linked debugging, visual plot previews
        inside the editor, and bit-exact runtime parity with \tv{} are non-goals of
        the open \lsp{} tier~\citep{pynescript2026}.
  \item \textbf{Security of metadata encryption.} Obscurity only; a determined party
        can extract the catalog from a binary.
  \item \textbf{Evaluation-in-IDE.} Richer offline mock-data evaluation panels and
        authenticated Pro preview commands are natural extensions.
\end{itemize}

% =============================================================================
\section{Conclusion}
\label{sec:conclusion}
% =============================================================================

Financial \textsc{Dsl}s need not remain trapped in browser-only editors. This paper
presented \texttt{pynescript-lsp}, a \texttt{pygls}-based Language Server for \pine{}
within the \pyne{} toolchain. By centering a shared \ast{}, caching parse/lint in a
workspace, and layering metadata-driven completion and hover over a generated builtin
catalog (\(\sim\)800+ items), the server delivers offline diagnostics, navigation,
semantic highlighting, inlay hints, and unparse-based formatting to any STDIO-capable
editor. Multi-surface delivery---VS~Code extension, thin Neovim/Zed/Emacs clients,
\cli{}, and a clean boundary with evaluation/\api{} tiers---shows how language
intelligence and runtime services can co-evolve without forcing online dependence for
editing. We believe the architecture is reusable for other bar-oriented and trading
\textsc{Dsl}s seeking IDE-grade tooling under the \lsp{} umbrella.

\section*{Acknowledgments}
We thank contributors to the HOOX / \pyne{} open-source stack and users who reported
editor integration issues across platforms.

\section*{Trademark disclaimer}
\trademarknote{}

% =============================================================================
\section*{References}
\bibliographystyle{plainnat}
\bibliography{../common/bibliography}

% =============================================================================
\appendix
\section{Capability comparison (fair matrix)}
\label{app:matrix}
% =============================================================================

Table~\ref{tab:compare} compares capability classes without claiming bit-exact parity
with proprietary hosted editors. Hosted environments excel at live chart coupling;
\pyne{} excels at offline, multi-editor, and CI workflows.

\begin{table}[h]
\centering
\caption{Capability matrix: browser-hosted \tv{} editor vs.\ \pyne{} local tooling
(qualitative; not a product claim of equivalence).}
\label{tab:compare}
\small
\begin{tabular}{@{}p{3.6cm}cc@{}}
\toprule
Capability & Hosted browser editor & \pyne{} \lsp{} + \cli{} \\
\midrule
Offline editing & Limited / none & Yes \\
Multi-editor (Vim/Emacs/\ldots) & No & Yes (\lsp{}) \\
Syntax highlighting & Yes & Yes (TextMate + semantic) \\
Diagnostics on edit & Platform-specific & Parse + lint (push/pull) \\
Builtin completion / hover & Yes (online) & Yes (local metadata) \\
Go-to-definition / refs & Platform-specific & Same-document \\
Deterministic format (CI) & Limited & Yes (unparse + \cli{}) \\
Live chart / strategy UI & Strong & Via optional Pro \api{} / external \\
Headless CI parse/lint & Weak & Strong \\
Open toolchain / self-host & No & Yes (AGPL package) \\
\bottomrule
\end{tabular}
\end{table}

\section{Example configuration fragments}
\label{app:config}

\paragraph{Neovim (conceptual).}
\begin{lstlisting}[style=python,numbers=none,basicstyle=\ttfamily\scriptsize,language={}]
require('lspconfig').pynescript.setup({})
-- server command: pynescript-lsp on PATH
\end{lstlisting}

\paragraph{Helix \texttt{languages.toml} (conceptual).}
\begin{lstlisting}[style=python,numbers=none,basicstyle=\ttfamily\scriptsize,language={}]
[[language]]
name = "pinescript"
file-types = ["pyne", "pine", "pinev5", "pinev6"]
command = "pynescript-lsp"
args = ["--stdio"]
\end{lstlisting}

\section{Software availability}
\label{app:avail}

Source: project repository associated with \pyne{} / HOOX~\citep{pynescript2026,hoox2026}.
Install language server: \texttt{pip install "hoox-pyne[lsp]"}. Launch:
\texttt{pynescript-lsp} or \texttt{python -m pynescript.langserver}.

\end{document}
